\documentclass{article}
\usepackage{graphicx} % Required for inserting images
\usepackage{amsmath} % For align* environment

\usepackage{algorithm}
\usepackage{algpseudocode}

\usepackage{amsmath, amsthm, amssymb}

\newtheorem{theorem}{Theorem}
\newtheorem{lemma}{Lemma}

\title{test}
\author{Joshua Roberge}
\date{October 2025}
\setlength{\parindent}{0pt}

\begin{document}

\begin{algorithm}
\caption{Right Rotate}
\begin{algorithmic}[1]
\State $\triangleright$ \text{Trigger by, left inbalance $B_F(3) = 2$ and 2's $B_F \geq 0$}
\Function{RightRotate}{cur}
    \State $\triangleright$ \text{Current = 3, we grab it left child (2) }
    \State cur\_left = cur.left
    \State $\triangleright$ \text{This is NULL in the given example}
    \State cur\_left\_right = cur\_left.right
    \State
    \State $\triangleright$ \text{Node 2 becomes root}
    \State cur\_left.right = cur
    \State $\triangleright$ \text{re-attach sub-tree}
    \State cur.left = cur\_left\_right
    \State
    \State $\triangleright$ \text{Recalculate Heights}
    \State cur.height = 1 
        + max(get\_avl\_height(cur.left),\\
        \hspace{7.5em} get\_avl\_height(cur.right))
    \State
    \State cur\_left.height = 1 
        + max(get\_avl\_height(cur\_left.left),\\
        \hspace{9em} get\_avl\_height(cur\_left.right))
    \State
    \State $\triangleright$ \text{Return the new root}
    \State \Return cur\_left
\EndFunction
\end{algorithmic}
\end{algorithm}


\begin{algorithm}
\caption{Left Rotate}
\begin{algorithmic}[1]
\State $\triangleright$ \text{Trigger by, right imbalance $B_F(1) = -2$ and 2's $B_F \leq 0$}
\Function{LeftRotate}{cur}
    \State $\triangleright$ \text{Current = 1, we grab its right child (2)}
    \State cur\_right = cur.right
    \State $\triangleright$ \text{This is NULL in the given example}
    \State cur\_right\_left = cur\_right.left
    \State
    \State $\triangleright$ \text{Node 2 becomes root}
    \State cur\_right.left = cur
    \State $\triangleright$ \text{re-attach sub-tree}
    \State cur.right = cur\_right\_left
    \State
    \State $\triangleright$ \text{Recalculate Heights}
    \State cur.height = 1 
        + max(get\_avl\_height(cur.left),\\
        \hspace{7.5em} get\_avl\_height(cur.right))
    \State
    \State cur\_right.height = 1 
        + max(get\_avl\_height(cur\_right.left),\\
        \hspace{9em} get\_avl\_height(cur\_right.right))
    \State
    \State $\triangleright$ \text{Return the new root}
    \State \Return cur\_right
\EndFunction
\end{algorithmic}
\end{algorithm}

\begin{algorithm}
\caption{Left-Right Rotate}
\begin{algorithmic}[1]
\State $\triangleright$ \text{Trigger by left-right imbalance ie}
\State \hspace{1.5em} \text{$B_F(3)>1$ and$B_F(1) < 0$}
\Function{Left-Right-Rotate}{cur}
    \State $\triangleright$ \text{Fixes the Zig-Zag,}
    \State \hspace{1.5em} \text{flattens the left imbalance}
    \State cur.left = LeftRotate(cur.left)
    \State $\triangleright$ \text{Perform right rotation,}
    \State \hspace{1.5em} \text{return new root (2)}
    \State \Return RightRotate(cur)
\EndFunction
\end{algorithmic}
\end{algorithm}

\begin{algorithm}
\caption{Right-Left Rotate}
\begin{algorithmic}[1]
\State $\triangleright$ \text{Trigger by right-left imbalance ie}
\State \hspace{1.5em} \text{$B_F(1)<-11$ and$B_F(3) > 0$}
\Function{Right-Left-Rotate}{cur}
    \State $\triangleright$ \text{Fixes the Zig-Zag,}
    \State \hspace{1.5em} \text{flattens the right imbalance}
    \State cur.right = RightRotate(cur.right)
    \State $\triangleright$ \text{Perform left rotation,}
    \State \hspace{1.5em} \text{return new root (2)}
    \State \Return LeftRotate(cur)
\EndFunction
\end{algorithmic}
\end{algorithm}





\setlength{\algorithmicindent}{.5em}

\begin{algorithm}
\caption{AVL-Insert}
\begin{algorithmic}[1]
\Function{Insert-AVL}{cur, value}
    \If{cur = NULL}
        \State $\triangleright$ \text{Creates a new Node, also our base case which end recursive calls}
        \State \Return CreateNode(value)
    \EndIf

    \State $\triangleright$  \text{Search Phase: recursive calls, will run until leaf is reach $log(n)$}
    \If{$value < cur.value$}
        \State cur.left = Insert-AVL(cur.left, value)
    \ElsIf{$value > cur.value$}
        \State cur.right = Insert-AVL(cur.right, value)
    \Else
        \State $\triangleright$ \text{Skip duplicates}
        \State \Return cur
    \EndIf

    \State $\triangleright$ \text{Rebalance/Metric Phase: these calls happen as we move back}\\ 
    \hspace{.7cm}\text{ up the tree after recursing downward $log(n)$}
    \State cur.height $\gets 1 + \max(\,Height(cur.left),\ Height(cur.right)\,)$
    
    \State cur $\gets$ Rebalance(cur)

    \State \Return cur

\EndFunction
\end{algorithmic}
\end{algorithm}




\setlength{\algorithmicindent}{.5em}

\begin{algorithm}
\caption{Search-AVL}
\begin{algorithmic}[1]
\State $\triangleright$ \text{A simple BST search, but we are guaranteed}\\
\hspace{.7cm}\text { to traverse a tree that is $h\le log(n)$}
\Function{Search-AVL}{root, value}
    \State cur = root
    \While{cur $\neq$ NULL} \quad{$\triangleright$ {while loop = $O(log(n))$}
        \If{$cur.value = value$}
            \State \Return cur
        \ElsIf{$value < cur.value$}
            \State cur = cur.left
        \Else
            \State cur = cur.right
        \EndIf
    \EndWhile
    \State \Return NULL
\EndFunction
\end{algorithmic}
\end{algorithm}





\setlength{\algorithmicindent}{.5em}

\begin{algorithm}
\caption{Delete-AVL}
\begin{algorithmic}[1]
\Function{Delete-AVL}{cur, value}
    \State $\triangleright$ \text{No Node found, base case}
    \If{cur == NULL}
        \State \Return cur
    \EndIf

    \State $\triangleright$ \text{Phase 1: search for node to delete $log(n)$}
    \If{$value < cur.value$}
        \State cur.left = Delete-AVL(cur.left, value)
    \ElsIf{$value > cur.value$}
        \State cur.right = Delete-AVL(cur.right, value)
    \Else
        \State $\triangleright$ \text{Phase 2: we found the Node, delete it $log(n)$}
        \If{$cur.left == NULL$}
            \State temp = cur.right
            \State delete(cur)
            \State \Return temp
        \ElsIf{cur.right == NULL}
            \State temp = cur.left
            \State delete(cur)
            \State \Return temp
        \Else
            \State $\triangleright$ \text{get min node can run until $log(n)$}
            \State min\_node = Get-Min-Node(cur.right)
            \State cur.value = min\_node.value
            \State cur.right = Delete-AVL(cur.right, min\_node.value)
        \EndIf
    \EndIf

    \State $\triangleright$ \text{Phase 3: move back up tree and rebalance + height metrics (log(n))}
    \State cur.height $= 1 + \max(\,Height(cur.left),\ Height(cur.right)\,)$
    \State cur = Rebalance(cur)
    \State \Return cur
\EndFunction
\end{algorithmic}
\end{algorithm}








\newpage
\clearpage

\begin{proof}[Direct Proof:]
Let $N_h$ be the minimum number of nodes in an AVL tree of height $h$. We claim that
\[
\forall\, h \ge 0,\quad h_n \le \mathcal{O}(\log(n)).
\]
To prove this claim, we first define the recursive relationship, then derive the upper bound of $N_h$, and finally finish by showing the's upper bound of $h$. When taken together, these steps will show that $ h_n \le \mathcal{O}(\log(n))$\\

\textbf{Recursive Relationship}:
\begin{quote}
 Let $N_L$ and $N_R$ be the minimum number of nodes for $N_h$'s left and right sub-tree. $N_h$ can be defined as a sum of $N_h = N_{L} + N_{R} +1$ (+1 is for the connecting node). Since we are dealing with the worst-case scenario and  we are constrained by our AVL's balance factor ($|B_F| \leq 1$), we derive the following relationship:
 \[
 N_L = N_{h-1} \text{, }
 N_R = N_{h-2}
 \]
 therefore,
 \[
  N_h = N_{h-1} + N_{h-2} + 1
 \] 
\end{quote}
\begin{proof}[\\Proof by Strong Induction (Proving $N_h$'s Lower Bound):]
\begin{quote}

\textbf{\\Base Cases:}
\begin{itemize}
    \item $N_1 = 1$, One root node
    \item $N_2 = 2 \ge 2^{\frac{2}{2}} $: One root node and child
    \item $N_3 = N_1 + N_2 + 2 = 2+1+1=4 \ge 2^{\frac{2}{2}} $: using the defined recursive relationship.
\end{itemize}

\textbf{Inductive Hypothesis:} Assume $k \ge 3$ and $\forall h$ s.t. $2 \ge h \le k$ that $N_h \ge 2^{\frac{h}{2}}$\\



\textbf{Inductive Step:}
\begin{align*}
N_{k+1} &= N_k + N_{k-1} + 1\\
        &> N_{k-1} + N_{k-1} + 1
        \quad\text{(Since we are proving the upper bound we can drop $N_k$ and $1$)}\\
        &> 2 \cdot N_{k-1}\\
        &> 2 \cdot 2^{(k-1)/2}
        \quad\text{(Inductive Hypothesis)}\\
        &> 2^{1 + (k-1)/2}\\
        &> 2^{(k+1)/2}
\end{align*}
Therefore: $N_{k+1} \ge 2^{\frac{k+1}{2}}$\\
\textbf{Conclusion}: $N_h \ge 2^{\frac{h}{2}}$ $\forall h\ge 2$
\end{quote}
\end{proof}


\textbf{$h$'s Upper Bound}
\begin{quote}
    
Taking the relationship derived above, we can derive $H$'s upper bound:
\begin{align*}
    N_h &> 2^{h/2}\\
    \log_2(N_h) &> \frac{h}{2}\\
    h &< 2\log_2(N_h)
\end{align*}
\end{quote}

\textbf{Conclusion:}
Since $n \ge N_h$ and $h < 2\log_2(N_h)$, we have $h = \mathcal{O}(\log n)$. Therefore, the height of an AVL tree with $n$ nodes is $\mathcal{O}(\log n)$.
\end{proof}
\end{document}
